\documentclass[10pt,a4paper]{article}
\usepackage{amsmath,amssymb,amsthm,graphicx}
\usepackage{amsmath,amssymb,amsthm}
\usepackage{hyperref}
\newtheorem{theorem}{Theorem}[section]
\newtheorem{definition}[theorem]{Definition}
\newtheorem{lemma}[theorem]{Lemma}
\newtheorem{claim}[theorem]{Claim}
\newtheorem{corollary}[theorem]{Corollary}
\DeclareMathOperator*{\ExpVal}{\mathbb E}
\newcommand{\N}{\mathbb N}
\newcommand{\ie}{i.\,e.}
\newcommand{\expref}[2]{\hyperref[#2]{#1~\ref*{#2}}}
\newcommand{\expeqref}[2]{\hyperref[#2]{#1~(\ref*{#2})}}
\title{Black-box deterministic expectation estimation for constant-width ROBPs from hitting sets}
\author{Kuan Cheng, William M. Hoza}
\date{doi:10.4086/toc.2022.v018a021; CC BY 3.0 Unported}
\begin{document}
\maketitle
\noindent\textit{Editorial note.} Editable author source from officialToC SourceZIP under articleCCBY3.0, with full selectedSection3 proof and necessary introductory definitions. Old same-sourceTheorem2.1 is not selected, repeated or counted; its one analogous-statement reference is an accurate external originalTheorem2.1 location. Original author formulas/indices/widths/space bounds retained without mathematical correction or added hypotheses. Publisherclass/template dependencies replaced by minimal presentation macros; no input/include or figures needed. Original section/theorem numbering restored, explicit source-local correctness and rounding/amplification included. Bibliography publisher formatting macros expanded to ordinary names/links; no bibliography package required. All production/validation exclusivelyYupeng; root owns compilation.\par\medskip
\noindent\textit{Allowed standard prerequisites (editorial).}\begin{itemize}
\item Read-once branching programs, hitting sets, oracle space models and standard hitting-set construction. OriginalSections1.1/1.3/1.5 definitions and fullSection3; publishedHozaZuckerman result cited only for uncounted unconditional instance
\end{itemize}\medskip
\setcounter{section}{1}
\section*{Model and definitions (original Section1)}
A width-$w$ length-$n$ read-once branching program (ROBP) is a directed graph consisting of $n+1$ layers with $w$ vertices per layer. There is a designated ``start vertex'' $v_{\text{start}}$ in the first layer. Every vertex not in the last layer has two outgoing edges labeled 0 and 1 leading to the next layer. An $n$-bit input string naturally identifies a path through the graph by reading from left to right. The program accepts or rejects this string depending on whether the path ends at the designated ``accept vertex'' $v_{\text{acc}}$ in the last layer.
\setcounter{theorem}{1}
\begin{definition}[Hitting set]
Let $\mathcal F$ be a class of functions $f:\{0,1\}^n\to\{0,1\}$. An $\epsilon$-hitting set for $\mathcal F$ is a set $H\subseteq\{0,1\}^n$ such that for every $f\in\mathcal F$ with $\ExpVal[f]\ge\epsilon$, there is some $x\in H$ such that $f(x)=1$.
\end{definition}
When we say ``a log-space hitting set,'' we mean a family of hitting sets $H_n\subseteq\{0,1\}^n$ such that given input $n\in\N$, the set $H_n$ can be enumerated in space $O(\log n)$. For such a family, $|H_n|\le\operatorname{poly}(n)$.
\setcounter{theorem}{4}
\begin{definition}[Deterministic sampler]
Let $\mathcal F$ be a class of functions $f:\{0,1\}^n\to\{0,1\}$. A deterministic $\epsilon$-sampler for $\mathcal F$ is a deterministic oracle algorithm $A$ that outputs a real number such that for every $f\in\mathcal F$, $|A^f-\ExpVal[f]|\le\epsilon$.
\end{definition}
Let $U_n$ denote the uniform distribution over $\{0,1\}^n$. For two strings $x,y$, let $x\circ y$ denote the concatenation of $x$ with $y$. Suppose an ROBP $f$ is clear from context. If $u$ and $v$ are vertices, let $p_{u\to v}$ be the probability that a random walk starting at $u$ reaches $v$. We use the shorthand $p_{\to v}=p_{v_{\text{start}}\to v}$ and $p_{u\to}=p_{u\to v_{\text{acc}}}$. We use $V_i$ to denote the set of vertices in the $i$-th layer of the ROBP, where $i\in\{0,1,\ldots,n\}$.

\setcounter{theorem}{6}
\begin{theorem} \label{thm:const-width-main}
Assume that for every constant $w$, for all $n \in \N$, there is a $\frac{1}{2}$-hitting set for width-$w$ length-$n$ ROBPs that can be computed in space $O(\log n)$. Then for every constant $w$, for all $n \in \N$ and all $\epsilon > 0$, there is a deterministic $\epsilon$-sampler for width-$w$ length-$n$ ROBPs that runs in space $O(\log(n/\epsilon))$.
\end{theorem}

More generally, given a $\frac{1}{2}$-hitting set for width-$w'$ length-$n'$ ROBPs for suitable values $w' = O(w^2)$ and $n' = O(n^2/\epsilon)$, we show how to construct a deterministic $\epsilon$-sampler for width-$w$ length-$n$ ROBPs. If the hitting set has space complexity $s$, then our deterministic sampler has space complexity $O(s + w \log(n/\epsilon))$. See \expref{Theorem}{thm:const-width-reduction} and \expref{Corollary}{cor:const-width-main-generalization}.

\begin{theorem}[Unconditional sampler] \label{thm:hz-sampler}
For every constant $w$, for all $n \in \N$ and all $\epsilon > 0$, there is a deterministic $\epsilon$-sampler for width-$w$ length-$n$ ROBPs running in space $O(\log^2 n + \log(1/\epsilon))$.
\end{theorem}

\setcounter{section}{2}
\section{Deterministic samplers for constant-width ROBPs} \label{sec:constwidth}

In this section, we will show how to use hitting sets to construct deterministic samplers for constant-width ROBPs, thereby proving \expref{Theorem}{thm:const-width-main}. Most of the work will go toward establishing the following reduction, which is meaningful even for slightly super-constant width.

\begin{theorem} \label{thm:const-width-reduction}
Let $w, n \in \N$ and let $\epsilon > 0$. Assume there is an $(\frac{\epsilon}{2n})$-hitting set $H$ for width-$(\binom{w}{2} + 1)$ length-$n$ ROBPs computable in space $s$. Then there is a deterministic $\epsilon$-sampler for width-$w$ length-$n$ ROBPs that runs in space $O(s + w \log(n/\epsilon))$.
\end{theorem}

Like Theorem 2.1 of the original publication, \expref{Theorem}{thm:const-width-reduction} technically ought to be phrased in terms of families of ROBPs. We should also clarify the model of space-bounded oracle algorithm. We assume that the sampler has write-only access to a ``query tape'' where it can write down an $n$-bit query string (the query string does not count against the sampler's space complexity). The sampler can then enter a special ``query'' state, which returns the result of the query into the algorithm's state and clears the query tape. This natural model was perhaps first studied by Ladner and Lynch \cite{LL76}.

\subsection{Setting up the reduction}
Toward proving \expref{Theorem}{thm:const-width-reduction}, we begin by setting up some notation. For any ROBP $f$ and a string $x \in \{0, 1\}^{\leq n}$, let $v_f(x)$ be the vertex reached when $f$ reads $x$. Furthermore, define
$$
p_f(x) = \ExpVal[f(x \circ U_{n - |x|})]\,,
$$
\ie, $p_f(x) = p_{v_f(x) \to}$.

Now, let $H \subseteq \{0, 1\}^n$ be an $\epsilon_H$-hitting set for width-$(\binom{w}{2} + 1)$ ROBPs. For $i \leq n$, let $H_i$ be the set of $i$-bit prefixes of strings in $H$, \ie, $H_i = \{x_1 x_2 \dots x_i : x \in H\}$. One can verify that $H_i$ is an $\epsilon_H$-hitting set for width-$(\binom{w}{2} + 1)$ length-$i$ ROBPs.

Let $f$ be the width-$w$ ROBP to which we have oracle access. Let $\lambda$ denote the empty string. Our goal is to estimate $p_f(\lambda)$. For each $i \leq n$, define an equivalence relation $\sim$ on $\{0, 1\}^i$ by the rule
$$
x \sim y \iff \forall z \in H_{n - i}, f(x \circ z) = f(y \circ z)\,.
$$
\begin{lemma} \label{lem:equiv}
If $x \sim y$, then $|p_f(x) - p_f(y)| < \epsilon_H$.
\end{lemma}

\begin{proof}
Let $i = |x| = |y|$. Define $g \colon \{0, 1\}^{n - i} \to \{0, 1\}$ by
$$
g(z) = f(x \circ z) \oplus f(y \circ z)\,.
$$
The function $g(z)$ can be computed by an ROBP of width $\binom{w}{2} + 1$: we have one state in $g$ for each unordered pair of states in $f$ to run the computations $f(x \circ z), f(y \circ z)$ in parallel, along with one additional $\bot$ state in $g$ to indicate that the two computations converged to the same state. If $|p_f(x) - p_f(y)| \geq \epsilon_H$, then $H_{n - i}$ hits $g$, hence $x \not \sim y$.
\end{proof}

Let $[x]$ denote the equivalence class of $x$, so $[x] \subseteq \{0, 1\}^{|x|}$. Our deterministic sampler will be based on numbers $\widetilde{p}_f([x]) \in [0, 1]$ for each equivalence class $[x]$. The definition of $\widetilde{p}_f$ will ensure that $\widetilde{p}_f([x]) \approx p_f(x)$ for typical values of $x$, although there might be some anomalous values of $x$ where $\widetilde{p}_f([x]) \not \approx p_f(x)$.

The definition of $\widetilde{p}_f([x])$ is inductive. For the base case, when $x \in \{0, 1\}^n$, define $\widetilde{p}_f([x]) = f(x)$. This is well-defined, because $x \sim y \implies f(x) = f(y)$. For the inductive step, suppose $x \in \{0, 1\}^i$ with $i < n$. Define
\begin{equation} \label{eqn:tilde-p-def}
\widetilde{p}_f([x]) = \max_{x' \in H_i \cap [x]} \left(\frac{1}{2} \widetilde{p}_f([x' \circ 0]) + \frac{1}{2} \widetilde{p}_f([x' \circ 1])\right)\,,
\end{equation}
with the convention that $\widetilde{p}_f([x]) = 0$ if $H_i \cap [x] = \varnothing$. Our sampler will output\footnote{Actually the sampler's output differs slightly from $\widetilde{p}_f([\lambda])$ due to rounding errors.} $\widetilde{p}_f([\lambda])$. (In \expref{Section}{sec:const-width-efficiency}, we will explain in more detail how to compute $\widetilde{p}_f([\lambda])$ in a space-efficient manner.)

\subsection{Correctness}

The upper bound on $\widetilde{p}_f([x])$ is straightforward:
\begin{claim} \label{clm:const-width-ub}
For every $i$, for every $x \in \{0, 1\}^{n - i}$,
$$
\widetilde{p}_f([x]) \leq p_f(x) + i \epsilon_H\,.
$$
\end{claim}

\begin{proof}
We proceed by induction on $i$. In the base case $i = 0$, $\widetilde{p}_f([x]) = f(x) = p_f(x)$. For the inductive step $i > 0$, we consider two cases. If $H_i \cap [x] = \varnothing$, then $\widetilde{p}_f([x]) = 0$ and the claim is trivial. Otherwise, there is some $x' \in H_i \cap [x]$ such that
\begin{align*}
\widetilde{p}_f([x]) &= \frac{1}{2} \widetilde{p}_f([x' \circ 0]) + \frac{1}{2} \widetilde{p}_f([x' \circ 1]) && (\text{\expeqref{Equation}{eqn:tilde-p-def}}) \\
&\leq \frac{1}{2} p_f(x' \circ 0) + \frac{1}{2} p_f(x' \circ 1) + (i - 1) \epsilon_H && \text{(Induction)} \\
&= p_f(x') + (i - 1) \epsilon_H \\
&< p_f(x) + i \epsilon_H && \text{(\expref{Lemma}{lem:equiv}.)} \qedhere
\end{align*}
\end{proof}

The lower bound is a little more subtle. If $u$ is a vertex in layer $i$ of $f$, we say that $u$ is \emph{$H$-reachable} if there is some $x \in H_i$ with $v_f(x) = u$. Otherwise, we say that $u$ is \emph{$H$-unreachable}. Let $\widetilde{f}$ be a width-$(w + 1)$ ROBP obtained from $f$ by replacing all $H$-unreachable nodes with reject nodes.\footnote{More precisely, add an extra node to each layer labeled $\bot$. In layers prior to the final layer, both outgoing edges from $\bot$ lead to $\bot$, and both outgoing edges from $H$-unreachable nodes lead to $\bot$. In the final layer, $H$-unreachable nodes and $\bot$ are reject nodes.}

\begin{claim} \label{clm:const-width-lb}
For every $i$, for every $x \in \{0, 1\}^{n - i}$,
$$
\widetilde{p}_f([x]) \geq p_{\widetilde{f}}(x)\,.
$$
\end{claim}

\begin{proof}
We proceed by induction on $i$. In the base case $i = 0$, $\widetilde{p}_f([x]) = f(x) \geq \widetilde{f}(x) = p_{\widetilde{f}}(x)$. For the inductive step $i > 0$, we consider two cases. If $f$ visits some $H$-unreachable node when it reads $x$, then $p_{\widetilde{f}}(x) = 0$ and the claim is trivial. Therefore, assume that when $f$ reads $x$, every node visited is $H$-reachable. Then there is some $x' \in H_{n - i}$ such that $v_f(x) = v_f(x')$. Of course when $f$ reads $x'$, every node visited is $H$-reachable, so
$$
v_{\widetilde{f}}(x') = v_f(x') = v_f(x) = v_{\widetilde{f}}(x)\,.
$$
Therefore,
\begin{align*}
p_{\widetilde{f}}(x) = p_{\widetilde{f}}(x') &= \frac{1}{2} p_{\widetilde{f}}(x' \circ 0) + \frac{1}{2} p_{\widetilde{f}}(x' \circ 1) \\
&\leq \frac{1}{2} \widetilde{p}_f([x' \circ 0]) + \frac{1}{2} \widetilde{p}_f([x' \circ 1]) && \text{(Induction)} \\
&\leq \widetilde{p}_f(x) && (\text{\expeqref{Equation}{eqn:tilde-p-def}.})
\end{align*}
(The last inequality uses the fact that $v_f(x) = v_f(x')$ and hence $x \sim x'$.)
\end{proof}

\begin{corollary} \label{cor:tilde-p-error}
$|\widetilde{p}_f([\lambda]) - \ExpVal[f]| \leq n \cdot \epsilon_H$.
\end{corollary}

\begin{proof}
By \expref{Claim}{clm:const-width-ub},
$$
\widetilde{p}_f([\lambda]) \leq p_f(\lambda) + n \cdot \epsilon_H = \ExpVal[f] + n \cdot \epsilon_H\,.
$$
In the other direction, by \expref{Claim}{clm:const-width-lb},
$$
\widetilde{p}_f([\lambda]) \geq p_{\widetilde{f}}(\lambda) = \ExpVal[\widetilde{f}]\,.
$$
Define $g \colon \{0, 1\}^n \to \{0, 1\}$ by
$$
g(x) = 1 \iff \text{when $f$ reads $x$, an $H$-unreachable node is visited}.
$$
Then $g$ can be computed by a width-$(w + 1)$ ROBP by a construction very similar to that of $\widetilde{f}$. By construction, $g$ rejects every string in $H$. Therefore, $\ExpVal[g] < \epsilon_H$. Furthermore, $g(x) = 0 \implies f(x) = \widetilde{f}(x)$. Therefore, $|\ExpVal[\widetilde{f}] - \ExpVal[f]| < \epsilon_H$, so $\widetilde{p}_f([\lambda]) > \ExpVal[f] - \epsilon_H$.
\end{proof}


\subsection{Efficiently computing $\widetilde{p}_f([\lambda])$} \label{sec:const-width-efficiency}

To complete the proof of \expref{Theorem}{thm:const-width-reduction}, we just need to show how to efficiently compute $\widetilde{p}_f([\lambda])$. This is fairly straightforward from the definitions; the details follow.

\begin{proof}[Proof of \expref{Theorem}{thm:const-width-reduction}]
Say a string $x \in H_i$ is a \emph{representative} if it is the lexicographically first element of $[x] \cap H_i$. Let $x^{(i, 1)}, x^{(i, 2)}, \dots$ be an enumeration of the representatives in $H_i$ in lexicographic order. Given $i$, $j$, and oracle access to $f$, one can compute $x^{(i, j)}$ in space $O(s)$.

Our sampler works its way backward through the branching program, starting at layer $n$ and ending with layer $0$. The sampler stores data about layer $i$ and uses it when processing layer $i - 1$. Specifically, the data stored regarding layer $i$ consists of a list of numbers $p_{i, 1}, p_{i, 2}, \dots$, with the interpretation $p_{i, j} = \widetilde{p}_f([x^{(i, j)}])$, or rather $p_{i, j} \approx \widetilde{p}_f([x^{(i, j)}])$ due to rounding error.

For layer $n$, we can compute this value exactly by setting $p_{i, j} = f(x^{(i, j)})$. 
Given these values for layer $i + 1$, define $\overline{p}_f(x)$ for each $x \in \{0, 1\}^{i + 1}$ by
\[
\overline{p}_f(x) = \begin{cases}
p_{i + 1, j} & \text{ if $j$ is such that $x^{(i + 1, j)} \sim x$} \\
0 & \text{ if no such $j$ exists.}
\end{cases}
\]
(Note that $j$ is unique if it exists, and note that $\widetilde{p}_f([x]) = 0$ if such a $j$ does not exist.) Then, we compute each value $p_{i, j}$ by the rule
\begin{equation} \label{eqn:pij}
p_{i, j} := \max_{x' \in H_i \cap [x^{(i, j)}]} \left(\frac{1}{2} \overline{p}_f(x' \circ 0) + \frac{1}{2} \overline{p}_f(x' \circ 1)\right)\,,
\end{equation}
with the convention that $p_{i, j} = 0$ if $H_i \cap [x^{(i, j)}] = \varnothing$.

The sampler performs the arithmetic in \expeqref{Equation}{eqn:pij} to within $\lceil \log(2n/\epsilon)\rceil$ bits of precision. This ensures that the rounding error is not too large in each step; by induction, $|p_{i, j} - \widetilde{p}_f([x^{(i, j)}])| \leq \frac{\epsilon (n - i)}{2n}$. The sampler outputs $p_{0, 1}$, which is within $\epsilon$ of $\ExpVal[f]$ by \expref{Corollary}{cor:tilde-p-error}, since $\epsilon_H = \frac{\epsilon}{2n}$.

The number of vertices in each layer of $f$ is at most $w$, so the number of equivalence classes in $\{0, 1\}^i$ is also at most $w$. Therefore, there are at most $w$ representatives in $H_i$, and hence there are only $w$ numbers $p_{i, j}$ being stored for each layer. Storing those numbers for the layer currently being processed and the layer most recently processed takes $O(w \log(n/\epsilon))$ bits of space, so overall, the space complexity of the sampler is $O(s + w \log(n/\epsilon))$ as claimed.
\end{proof}

Interestingly, the sampler of \expref{Theorem}{thm:const-width-reduction} can be implemented to be non-adaptive, because it only queries $f$ at strings of the form $x \circ y$ or $x \circ b \circ z$, where $x \in H_i$, $y \in H_{n - i}$, $b \in \{0, 1\}$, and $z \in H_{n - i - 1}$.

\subsection{Applying the reduction}

\begin{proof}[Proof of \expref{Theorem}{thm:hz-sampler}]
Hoza and Zuckerman constructed an $\left(\frac{\epsilon}{2n}\right)$-hitting set $H$ even for \emph{polynomial}-width ROBPs that can be computed in space $O(\log^2 n + \log(1/\epsilon))$ \cite{HZ20}. Combining this result with \expref{Theorem}{thm:const-width-reduction} immediately proves \expref{Theorem}{thm:hz-sampler}.
\end{proof}

To prove \expref{Theorem}{thm:const-width-main}, we must first amplify the assumed $\frac{1}{2}$-hitting set to get an $\left(\frac{\epsilon}{2n}\right)$-hitting set. This is straightforward, although we must pay a small penalty in terms of width, length, and cardinality.
\begin{lemma} \label{lem:hitting-set-amplification}
Suppose $H$ is a $\frac{1}{2}$-hitting set for width-$(w + 1)$ length-$(nm)$ ROBPs. Divide each string $x \in H$ into blocks of length $n$, $x = x^{(1)} \circ x^{(2)} \circ \cdots \circ x^{(m)}$. Let $H' = \{x^{(i)} : x \in H, i \in [m]\}$. Then $H'$ is a $(\frac{1}{m})$-hitting set for width-$w$ length-$n$ ROBPs.
\end{lemma}

\begin{proof}
Let $f$ be a width-$w$ length-$n$ ROBP with $\ExpVal[f] \geq 1/m$. Define $g \colon (\{0, 1\}^n)^m \to \{0, 1\}$ by
$$
g(x^{(1)} \circ \dots \circ x^{(m)}) = \bigvee_{i \in [m]} f(x^{(i)})\,.
$$
Then $g$ can be computed by a width-$(w + 1)$ ROBP. Furthermore,
$$
\ExpVal[g] = 1 - (1 - \ExpVal[f])^m \geq 1 - \left(1 - \frac{1}{m}\right)^m > \frac{1}{2}\,.
$$
Therefore, $H$ hits $g$, hence $H'$ hits $f$.
\end{proof}

Combining~\expref{Lemma}{lem:hitting-set-amplification} with \expref{Theorem}{thm:const-width-reduction} proves the following corollary, which immediately implies \expref{Theorem}{thm:const-width-main}.

\begin{corollary} \label{cor:const-width-main-generalization}
Let $w, n \in \N$ and let $\epsilon > 0$. Let $w' = \binom{w}{2} + 2$ and $n' = n \cdot \lceil 2n / \epsilon \rceil$. Assume there is a $\frac{1}{2}$-hitting set $H$ for width-$w'$ length-$n'$ ROBPs computable in space $s$. Then there is a deterministic $\epsilon$-sampler for width-$w$ length-$n$ ROBPs that runs in space $O(s + w \log(n/\epsilon))$.
\end{corollary}

\begin{thebibliography}{29}
\bibitem{HZ20}
{\sc William~M. Hoza and David Zuckerman}: Simple optimal hitting sets for
  small-success {$\mathbf{RL}$}.
\newblock {\em SIAM J. Comput.}, 49(4):811--820, 2020.
\newblock Preliminary version in \href{https://doi.org/10.1109/FOCS.2018.00015}{FOCS'18}.
\newblock [\href{https://doi.org/10.1137/19M1268707}{doi:10.1137/19M1268707}, ECCC TR18-063]

\bibitem{LL76}
{\sc Richard~E. Ladner and Nancy~A. Lynch}: Relativization of questions about
  log space computability.
\newblock {\em Math. Sys. Theory}, 10(1):19--32, 1976.
\newblock [\href{https://doi.org/10.1007/BF01683260}{doi:10.1007/BF01683260}]

\end{thebibliography}
\end{document}
